% Options for packages loaded elsewhere
\PassOptionsToPackage{unicode}{hyperref}
\PassOptionsToPackage{hyphens}{url}
\PassOptionsToPackage{dvipsnames,svgnames,x11names}{xcolor}
\documentclass[
  11pt,
  a4paper,
]{article}
\usepackage{xcolor}
\usepackage[margin=1in]{geometry}
\usepackage{amsmath,amssymb}
\setcounter{secnumdepth}{-\maxdimen} % remove section numbering
\usepackage{iftex}
\ifPDFTeX
  \usepackage[T1]{fontenc}
  \usepackage[utf8]{inputenc}
  \usepackage{textcomp} % provide euro and other symbols
\else % if luatex or xetex
  \usepackage{unicode-math} % this also loads fontspec
  \defaultfontfeatures{Scale=MatchLowercase}
  \defaultfontfeatures[\rmfamily]{Ligatures=TeX,Scale=1}
\fi
\usepackage{lmodern}
\ifPDFTeX\else
  % xetex/luatex font selection
  \ifLuaTeX
    \usepackage{luaotfload}
    \directlua{luaotfload.add_fallback("mainfontfallback",{
      "NotoSerifBengali:mode=harf"
    })}
  \fi
  \setmainfont[,RawFeature={fallback=mainfontfallback}]{Latin Modern Roman}
  \ifLuaTeX
    \usepackage{luaotfload}
    \directlua{luaotfload.add_fallback("monofontfallback",{
      "NotoSansBengali:mode=harf"
    })}
  \fi
  \setmonofont[,RawFeature={fallback=monofontfallback}]{DejaVu Sans Mono}
\fi
% Use upquote if available, for straight quotes in verbatim environments
\IfFileExists{upquote.sty}{\usepackage{upquote}}{}
\IfFileExists{microtype.sty}{% use microtype if available
  \usepackage[]{microtype}
  \UseMicrotypeSet[protrusion]{basicmath} % disable protrusion for tt fonts
}{}
\makeatletter
\@ifundefined{KOMAClassName}{% if non-KOMA class
  \IfFileExists{parskip.sty}{%
    \usepackage{parskip}
  }{% else
    \setlength{\parindent}{0pt}
    \setlength{\parskip}{6pt plus 2pt minus 1pt}}
}{% if KOMA class
  \KOMAoptions{parskip=half}}
\makeatother
\usepackage{color}
\usepackage{fancyvrb}
\newcommand{\VerbBar}{|}
\newcommand{\VERB}{\Verb[commandchars=\\\{\}]}
\DefineVerbatimEnvironment{Highlighting}{Verbatim}{commandchars=\\\{\}}
% Add ',fontsize=\small' for more characters per line
\usepackage{framed}
\definecolor{shadecolor}{RGB}{248,248,248}
\newenvironment{Shaded}{\begin{snugshade}}{\end{snugshade}}
\newcommand{\AlertTok}[1]{\textcolor[rgb]{0.94,0.16,0.16}{#1}}
\newcommand{\AnnotationTok}[1]{\textcolor[rgb]{0.56,0.35,0.01}{\textbf{\textit{#1}}}}
\newcommand{\AttributeTok}[1]{\textcolor[rgb]{0.13,0.29,0.53}{#1}}
\newcommand{\BaseNTok}[1]{\textcolor[rgb]{0.00,0.00,0.81}{#1}}
\newcommand{\BuiltInTok}[1]{#1}
\newcommand{\CharTok}[1]{\textcolor[rgb]{0.31,0.60,0.02}{#1}}
\newcommand{\CommentTok}[1]{\textcolor[rgb]{0.56,0.35,0.01}{\textit{#1}}}
\newcommand{\CommentVarTok}[1]{\textcolor[rgb]{0.56,0.35,0.01}{\textbf{\textit{#1}}}}
\newcommand{\ConstantTok}[1]{\textcolor[rgb]{0.56,0.35,0.01}{#1}}
\newcommand{\ControlFlowTok}[1]{\textcolor[rgb]{0.13,0.29,0.53}{\textbf{#1}}}
\newcommand{\DataTypeTok}[1]{\textcolor[rgb]{0.13,0.29,0.53}{#1}}
\newcommand{\DecValTok}[1]{\textcolor[rgb]{0.00,0.00,0.81}{#1}}
\newcommand{\DocumentationTok}[1]{\textcolor[rgb]{0.56,0.35,0.01}{\textbf{\textit{#1}}}}
\newcommand{\ErrorTok}[1]{\textcolor[rgb]{0.64,0.00,0.00}{\textbf{#1}}}
\newcommand{\ExtensionTok}[1]{#1}
\newcommand{\FloatTok}[1]{\textcolor[rgb]{0.00,0.00,0.81}{#1}}
\newcommand{\FunctionTok}[1]{\textcolor[rgb]{0.13,0.29,0.53}{\textbf{#1}}}
\newcommand{\ImportTok}[1]{#1}
\newcommand{\InformationTok}[1]{\textcolor[rgb]{0.56,0.35,0.01}{\textbf{\textit{#1}}}}
\newcommand{\KeywordTok}[1]{\textcolor[rgb]{0.13,0.29,0.53}{\textbf{#1}}}
\newcommand{\NormalTok}[1]{#1}
\newcommand{\OperatorTok}[1]{\textcolor[rgb]{0.81,0.36,0.00}{\textbf{#1}}}
\newcommand{\OtherTok}[1]{\textcolor[rgb]{0.56,0.35,0.01}{#1}}
\newcommand{\PreprocessorTok}[1]{\textcolor[rgb]{0.56,0.35,0.01}{\textit{#1}}}
\newcommand{\RegionMarkerTok}[1]{#1}
\newcommand{\SpecialCharTok}[1]{\textcolor[rgb]{0.81,0.36,0.00}{\textbf{#1}}}
\newcommand{\SpecialStringTok}[1]{\textcolor[rgb]{0.31,0.60,0.02}{#1}}
\newcommand{\StringTok}[1]{\textcolor[rgb]{0.31,0.60,0.02}{#1}}
\newcommand{\VariableTok}[1]{\textcolor[rgb]{0.00,0.00,0.00}{#1}}
\newcommand{\VerbatimStringTok}[1]{\textcolor[rgb]{0.31,0.60,0.02}{#1}}
\newcommand{\WarningTok}[1]{\textcolor[rgb]{0.56,0.35,0.01}{\textbf{\textit{#1}}}}
\usepackage{longtable,booktabs,array}
\usepackage{calc} % for calculating minipage widths
% Correct order of tables after \paragraph or \subparagraph
\usepackage{etoolbox}
\makeatletter
\patchcmd\longtable{\par}{\if@noskipsec\mbox{}\fi\par}{}{}
\makeatother
% Allow footnotes in longtable head/foot
\IfFileExists{footnotehyper.sty}{\usepackage{footnotehyper}}{\usepackage{footnote}}
\makesavenoteenv{longtable}
\setlength{\emergencystretch}{3em} % prevent overfull lines
\providecommand{\tightlist}{%
  \setlength{\itemsep}{0pt}\setlength{\parskip}{0pt}}
% paper_preamble.tex — shared LaTeX preamble for the PCMA/SI paper set.
%
% Compiled with lualatex (NOT pdflatex) because several papers contain
% Bengali script: the CLL paper realises one cognitive state into
% English, Bengali and Banglish, so the Bengali output is EVIDENCE and
% cannot be transliterated away without destroying the point being made.
%
% Font selection is NOT done here. pandoc >= 3 has native
% `mainfontfallback` / `monofontfallback` variables that emit the
% luaotfload fallback registration in the correct position relative to
% unicode-math; doing it by hand in this header lands after unicode-math
% and breaks font loading outright. See tools/build_papers.sh.

\usepackage{amsmath}
\usepackage{amssymb}
\usepackage{mathtools}   % \xleftrightarrow etc. used in cll.md §5
\usepackage{booktabs}

% ---------------------------------------------------------------------
% Unicode symbols appearing in PROSE rather than inside $...$.
%
% The papers were written as markdown and read as markdown, so they use
% real characters — "the ω-boost", "α-child", "ΔR coherence" — instead of
% LaTeX macros. Under unicode-math those codepoints are math-mode-only
% and abort the build ("\symbballowed only in math mode"). Mapping each
% to \ensuremath keeps the markdown readable on GitHub AND compilable,
% which is the whole point of keeping one source.
% ---------------------------------------------------------------------
\usepackage{newunicodechar}

% Greek, lower case
\newunicodechar{α}{\ensuremath{\alpha}}
\newunicodechar{β}{\ensuremath{\beta}}
\newunicodechar{γ}{\ensuremath{\gamma}}
\newunicodechar{δ}{\ensuremath{\delta}}
\newunicodechar{ε}{\ensuremath{\varepsilon}}
\newunicodechar{η}{\ensuremath{\eta}}
\newunicodechar{θ}{\ensuremath{\theta}}
\newunicodechar{κ}{\ensuremath{\kappa}}
\newunicodechar{λ}{\ensuremath{\lambda}}
\newunicodechar{μ}{\ensuremath{\mu}}
\newunicodechar{π}{\ensuremath{\pi}}
\newunicodechar{ρ}{\ensuremath{\rho}}
\newunicodechar{σ}{\ensuremath{\sigma}}
\newunicodechar{τ}{\ensuremath{\tau}}
\newunicodechar{φ}{\ensuremath{\varphi}}
\newunicodechar{χ}{\ensuremath{\chi}}
\newunicodechar{ψ}{\ensuremath{\psi}}
\newunicodechar{ω}{\ensuremath{\omega}}
\newunicodechar{ν}{\ensuremath{\nu}}
\newunicodechar{ξ}{\ensuremath{\xi}}
\newunicodechar{ζ}{\ensuremath{\zeta}}
\newunicodechar{ι}{\ensuremath{\iota}}

% Greek, upper case
\newunicodechar{Δ}{\ensuremath{\Delta}}
\newunicodechar{Φ}{\ensuremath{\Phi}}
\newunicodechar{Σ}{\ensuremath{\Sigma}}
\newunicodechar{Ψ}{\ensuremath{\Psi}}
\newunicodechar{Ω}{\ensuremath{\Omega}}
\newunicodechar{Γ}{\ensuremath{\Gamma}}
\newunicodechar{Λ}{\ensuremath{\Lambda}}
\newunicodechar{Θ}{\ensuremath{\Theta}}
\newunicodechar{Π}{\ensuremath{\Pi}}

% Operators and relations
\newunicodechar{×}{\ensuremath{\times}}
\newunicodechar{·}{\ensuremath{\cdot}}
\newunicodechar{±}{\ensuremath{\pm}}
\newunicodechar{−}{\ensuremath{-}}
\newunicodechar{≈}{\ensuremath{\approx}}
\newunicodechar{≤}{\ensuremath{\leq}}
\newunicodechar{≥}{\ensuremath{\geq}}
\newunicodechar{≠}{\ensuremath{\neq}}
\newunicodechar{∈}{\ensuremath{\in}}
\newunicodechar{∂}{\ensuremath{\partial}}
\newunicodechar{∇}{\ensuremath{\nabla}}
\newunicodechar{√}{\ensuremath{\sqrt{}}}
\newunicodechar{∞}{\ensuremath{\infty}}
\newunicodechar{∝}{\ensuremath{\propto}}
\newunicodechar{∑}{\ensuremath{\sum}}
\newunicodechar{∫}{\ensuremath{\int}}
\newunicodechar{ℝ}{\ensuremath{\mathbb{R}}}
\newunicodechar{ℕ}{\ensuremath{\mathbb{N}}}
\newunicodechar{‖}{\ensuremath{\|}}
\newunicodechar{⊕}{\ensuremath{\oplus}}
\newunicodechar{⊗}{\ensuremath{\otimes}}
\newunicodechar{∀}{\ensuremath{\forall}}
\newunicodechar{∃}{\ensuremath{\exists}}
\newunicodechar{¬}{\ensuremath{\neg}}
\newunicodechar{∧}{\ensuremath{\wedge}}
\newunicodechar{∨}{\ensuremath{\vee}}

% Arrows
\newunicodechar{→}{\ensuremath{\rightarrow}}
\newunicodechar{←}{\ensuremath{\leftarrow}}
\newunicodechar{↔}{\ensuremath{\leftrightarrow}}
\newunicodechar{↑}{\ensuremath{\uparrow}}
\newunicodechar{↓}{\ensuremath{\downarrow}}
\newunicodechar{⇒}{\ensuremath{\Rightarrow}}
\newunicodechar{⇐}{\ensuremath{\Leftarrow}}
\newunicodechar{⇔}{\ensuremath{\Leftrightarrow}}

% Super/subscripts written as literal characters
\newunicodechar{²}{\ensuremath{^{2}}}
\newunicodechar{³}{\ensuremath{^{3}}}
\newunicodechar{⁴}{\ensuremath{^{4}}}
\newunicodechar{⁵}{\ensuremath{^{5}}}
\newunicodechar{⁶}{\ensuremath{^{6}}}
\newunicodechar{₀}{\ensuremath{_{0}}}
\newunicodechar{₁}{\ensuremath{_{1}}}
\newunicodechar{₂}{\ensuremath{_{2}}}
\newunicodechar{₃}{\ensuremath{_{3}}}
\newunicodechar{ᵀ}{\ensuremath{^{\mathsf{T}}}}

% Status glyphs. The papers use emoji as status markers in tables
% (implemented / failed / warning). Mapped to typographic equivalents:
% a colour emoji font in lualatex is unreliable, and the MEANING is what
% carries — a reader needs to see "this one failed", not a red cross.
\newunicodechar{✅}{\ensuremath{\checkmark}}
\newunicodechar{✓}{\ensuremath{\checkmark}}
\newunicodechar{❌}{\ensuremath{\times}}
\newunicodechar{✗}{\ensuremath{\times}}
\newunicodechar{⚠}{\textbf{!}}
\newunicodechar{🧩}{\textbf{[piece]}}
\newunicodechar{📋}{\textbf{[list]}}
\newunicodechar{🟡}{\textbf{[partial]}}
\newunicodechar{🔴}{\textbf{[dead]}}
\newunicodechar{🟢}{\textbf{[live]}}
\newunicodechar{▼}{\ensuremath{\blacktriangledown}}
\newunicodechar{►}{\ensuremath{\blacktriangleright}}
\newunicodechar{□}{\ensuremath{\square}}
\newunicodechar{█}{\rule{0.6em}{0.9ex}}
\newunicodechar{•}{\textbullet}
\newunicodechar{†}{\textdagger}
\newunicodechar{‑}{-}
\newunicodechar{‐}{-}
\newunicodechar{…}{\ldots}
\newunicodechar{।}{\textbar}   % danda: rendered by the Bengali fallback in text
\newunicodechar{⁻}{\ensuremath{^{-}}}

% Zero-width formatting codepoints. These are invisible in an editor but
% are real characters to LaTeX: an emoji written as "⚠️" is TWO
% codepoints (U+26A0 U+FE0F), so declaring only the visible one leaves a
% variation selector with no definition and the build dies on a symbol
% nobody can see. Both are mapped to nothing.
\newunicodechar{^^^^fe0f}{}    % VARIATION SELECTOR-16
\newunicodechar{^^^^200d}{}    % ZERO WIDTH JOINER
\newunicodechar{^^^^200b}{}    % ZERO WIDTH SPACE

\usepackage{longtable}
\usepackage{array}

% Comparison matrices and cost ledgers are wide. Let them shrink to the
% text block rather than overflow the margin.
\usepackage{adjustbox}

% Code blocks: wrap long lines instead of running off the page. Several
% papers quote real C++ with long identifiers.
\usepackage{fvextra}
\fvset{breaklines=true,breakanywhere=true,fontsize=\small}

\setlength{\parskip}{0.5em}
\setlength{\parindent}{0pt}

\usepackage{fancyhdr}
\pagestyle{fancy}
\fancyhf{}
\fancyfoot[C]{\thepage}
\fancyhead[L]{\small\itshape Synthetic Intelligence — research lineage}
\fancyhead[R]{\small\itshape irfan@furylogic.com}
\renewcommand{\headrulewidth}{0.2pt}
\setlength{\headheight}{14pt}
\setlength{\headsep}{18pt}
\usepackage{bookmark}
\IfFileExists{xurl.sty}{\usepackage{xurl}}{} % add URL line breaks if available
\urlstyle{same}
\hypersetup{
  pdftitle={Einstein\textquotesingle s Epistemology as a Running Architecture --- Induction, Refutation and Impossibility},
  pdfauthor={Irfan Mahir --- Independent Researcher, Dhaka, Bangladesh --- irfan@furylogic.com},
  colorlinks=true,
  linkcolor={black},
  filecolor={Maroon},
  citecolor={Blue},
  urlcolor={Blue},
  pdfcreator={LaTeX via pandoc}}

\title{Einstein\textquotesingle s Epistemology as a Running Architecture
--- Induction, Refutation and Impossibility}
\author{Irfan Mahir --- Independent Researcher, Dhaka, Bangladesh ---
irfan@furylogic.com}
\date{}

\begin{document}
\maketitle

{
\hypersetup{linkcolor=}
\setcounter{tocdepth}{2}
\tableofcontents
}
\section{Einstein's Epistemology as a Running Architecture ---
Induction, Refutation and
Impossibility}\label{einsteins-epistemology-as-a-running-architecture-induction-refutation-and-impossibility}

\textbf{Irfan Mahir}\\
\emph{Independent Researcher, Dhaka, Bangladesh}\\
irfan@furylogic.com

\begin{quote}
The J-A-S cycle, and the register that records the warrant of every
``cannot''.
\end{quote}

\begin{center}\rule{0.5\linewidth}{0.5pt}\end{center}

\subsection{Contents of This Volume}\label{contents-of-this-volume}

\begin{longtable}[]{@{}
  >{\raggedright\arraybackslash}p{(\linewidth - 4\tabcolsep) * \real{0.3333}}
  >{\raggedright\arraybackslash}p{(\linewidth - 4\tabcolsep) * \real{0.3333}}
  >{\raggedright\arraybackslash}p{(\linewidth - 4\tabcolsep) * \real{0.3333}}@{}}
\toprule\noalign{}
\begin{minipage}[b]{\linewidth}\raggedright
\#
\end{minipage} & \begin{minipage}[b]{\linewidth}\raggedright
Paper
\end{minipage} & \begin{minipage}[b]{\linewidth}\raggedright
File
\end{minipage} \\
\midrule\noalign{}
\endhead
\bottomrule\noalign{}
\endlastfoot
1 & JAS: Einstein's Epistemology as a Running Architecture ---
Induction, Refutation, and the Warrant of Impossibility &
\texttt{jas.md} \\
\end{longtable}

\begin{center}\rule{0.5\linewidth}{0.5pt}\end{center}

\newpage

\subsection{JAS: Einstein's Epistemology as a Running Architecture ---
Induction, Refutation, and the Warrant of
Impossibility}\label{jas-einsteins-epistemology-as-a-running-architecture-induction-refutation-and-the-warrant-of-impossibility}

\begin{quote}
\emph{``A system that cannot be wrong is not doing science. It is doing
bookkeeping with a confident tone.''}
\end{quote}

\begin{center}\rule{0.5\linewidth}{0.5pt}\end{center}

\subsubsection{Abstract}\label{abstract}

Learning systems are conventionally organised around an objective that
is minimised. This paper describes an architecture organised instead
around a claim that can be destroyed. We implement Einstein's account of
physical theory --- the \textbf{J-A-S cycle}, in which an unjustified
\emph{jump} (J) from sense experience produces axioms (A), from which
theorems (S) are deduced and tested back against experience --- as a
running loop with an external oracle. Induction is counting; deduction
is graph traversal; refutation is a compiler exit code. There is no loss
function, no gradient, no fitted parameter, and no policy.

We report a deliberately false induction, its experimental refutation,
and its revision into a restricted claim that survives on held-out
symbols. We then introduce the \textbf{impossibility register}, which
records the \emph{warrant} of every ``X cannot be done'' claim in the
system. The central argument is that warrant determines what a
counterexample means: refuting an INDUCED universal ends it, while a
counterexample to a THEOREM is a \textbf{scope error until proven
otherwise}. Two of this project's own impossibility theorems were
escaped by leaving their scope; neither was violated. We argue that this
distinction --- routinely collapsed in machine-learning discourse --- is
what separates a field that knows where its exits are from one that
either stalls at a proof it misread or celebrates a breakthrough it did
not achieve.

\begin{center}\rule{0.5\linewidth}{0.5pt}\end{center}

\subsubsection{1. The Diagram, and Why Its Asymmetry
Matters}\label{the-diagram-and-why-its-asymmetry-matters}

Einstein described the structure of physical theory not as a ladder from
observation to law, but as a cycle:

\begin{verbatim}
                  +---------------------------+
                  |        Axioms (A)         |
                  +---------------------------+
                     /                     \
                    v                       v
                +-----+                   +-----+
                |  S  |                   | S'  |   deduced theorems
                +-----+                   +-----+
                   |                         |
                   | experiment              | experiment
                   v                         v
~~~~~~~~~~~~~~~~~~~~~~~~~~~~~~~~~~~~~~~~~~~~~~~~~~~~
               Sense Experience (E)
~~~~~~~~~~~~~~~~~~~~~~~~~~~~~~~~~~~~~~~~~~~~~~~~~~~~
       \                                     /
        +------------- Jump (J) ------------+
\end{verbatim}

The load-bearing feature is that \textbf{J is not logical}. No inference
rule carries you from a set of observations to an axiom. The arrow is
inductive, unjustified, and revisable. Everything below the axioms is
deductive and mechanical; everything from experience up to the axioms is
a leap.

This asymmetry is what most learning architectures erase. Gradient
descent presents itself as a continuous ascent from data to model, with
no moment at which something unjustified is asserted. The claim of this
paper is that erasing the leap does not remove it --- it only removes
the system's ability to notice that it made one, and therefore its
ability to retract it.

\paragraph{1.1 What the substrate already had, and what was
missing}\label{what-the-substrate-already-had-and-what-was-missing}

Prior work in this lineage supplied every element of the diagram except
the leap:

\begin{longtable}[]{@{}
  >{\raggedright\arraybackslash}p{(\linewidth - 4\tabcolsep) * \real{0.3333}}
  >{\raggedright\arraybackslash}p{(\linewidth - 4\tabcolsep) * \real{0.3333}}
  >{\raggedright\arraybackslash}p{(\linewidth - 4\tabcolsep) * \real{0.3333}}@{}}
\toprule\noalign{}
\begin{minipage}[b]{\linewidth}\raggedright
Element
\end{minipage} & \begin{minipage}[b]{\linewidth}\raggedright
Mechanism
\end{minipage} & \begin{minipage}[b]{\linewidth}\raggedright
Where
\end{minipage} \\
\midrule\noalign{}
\endhead
\bottomrule\noalign{}
\endlastfoot
\textbf{E} --- sense experience & mechanical AST extraction from source
& \texttt{tools/extract\_ast\_axioms.py} \\
\textbf{A} --- axioms & typed concept graph & \texttt{data/*.axioms} \\
\textbf{S} --- theorems & forward chaining over nine primitives & 100
derived edges, premises cited \\
\textbf{experiment} & compile-verify against an external oracle &
\texttt{synth.cpp} \\
\textbf{J} --- the jump & \emph{absent} & --- \\
\end{longtable}

Without J the system could only elaborate what it was given. It could
derive \texttt{human\ is\_a\ great\_ape} transitively, and it could
discover by experiment that the extractor was wrong about a function's
arity --- but it could not \emph{assert a universal that nothing
entailed} and then stake it on an outcome that might destroy it.

\begin{center}\rule{0.5\linewidth}{0.5pt}\end{center}

\subsubsection{2. The Jump}\label{the-jump}

\paragraph{2.1 Reading the world, not being told about
it}\label{reading-the-world-not-being-told-about-it}

The system must first decide what counts as a function. It is not told.
It reads the graph: a node with an outgoing \texttt{calls} edge is a
function. That is sense experience --- a regularity observed in
structure that was mechanically extracted from source, with no human
labelling step.

\begin{verbatim}
functions  = { e.from : e ∈ E,  rel(e) = calls }
has_args   = { e.from : e ∈ E,  rel(e) = takes_arg }
\end{verbatim}

\paragraph{2.2 Leaping further than the
evidence}\label{leaping-further-than-the-evidence}

Counting gives an observed frequency. On the live graph, 55\% of nodes
identified as functions carry at least one \texttt{takes\_arg} edge.

The system then asserts:

\begin{quote}
\textbf{Induced axiom.} \emph{Every function takes at least one
argument.}
\end{quote}

This is deliberately stronger than the evidence. That gap is the entire
point. An induction that merely restated ``55\% of observed functions
have arguments'' would be a \textbf{summary}, and a summary cannot be
refuted --- it is already true. Only a universal can be destroyed by a
single case, and only something destructible is a scientific claim.

Provenance is recorded as \texttt{INDUCED}, the weakest warrant in the
system. Its sole justification is that nothing has yet destroyed it.

\begin{center}\rule{0.5\linewidth}{0.5pt}\end{center}

\subsubsection{3. Deduction and
Experiment}\label{deduction-and-experiment}

\paragraph{3.1 A -\textgreater{} S}\label{a---s}

From the induced axiom, a concrete and falsifiable prediction is deduced
for a specific symbol:

\begin{quote}
If every function takes at least one argument, then calling
\texttt{run()} with zero arguments must \textbf{fail to compile}.
\end{quote}

Deduction here is ordinary and mechanical. Nothing interesting happens
in this step, which is exactly right: the interesting steps are the leap
above it and the experiment below it.

\paragraph{3.2 S -\textgreater{} E: the oracle}\label{s---e-the-oracle}

The prediction is tested by emitting a translation unit and invoking the
C++ compiler:

\begin{Shaded}
\begin{Highlighting}[]
\PreprocessorTok{\#include }\ImportTok{"physics.hpp"}
\DataTypeTok{int}\NormalTok{ main}\OperatorTok{()\{}\NormalTok{ PhysicsEngine px}\OperatorTok{;}\NormalTok{ px}\OperatorTok{.}\NormalTok{run}\OperatorTok{();} \ControlFlowTok{return} \DecValTok{0}\OperatorTok{;} \OperatorTok{\}}
\end{Highlighting}
\end{Shaded}

The compiler is the correct oracle for three reasons. It is
\textbf{deterministic} --- the same input yields the same verdict. It is
\textbf{external} --- it knows nothing about the substrate, shares no
representation with it, and cannot be argued with. And it is
\textbf{indifferent} --- it has no stake in whether the axiom survives.

An exit code is the refutation.

\paragraph{3.3 The refutation}\label{the-refutation}

\begin{verbatim}
tick            predicted fails       observed compiles    -> REFUTES AXIOM
run             predicted fails       observed compiles    -> REFUTES AXIOM
reset_field     predicted fails       observed compiles    -> REFUTES AXIOM
drive           predicted fails       observed fails       -> CONSISTENT
\end{verbatim}

\texttt{tick()}, \texttt{run()} and \texttt{reset\_field()} compile with
zero arguments. The universal is false. It is retracted, and the
retraction is logged to \texttt{data/refuted.axioms} --- because a
failed induction is knowledge: it is a region of hypothesis space now
known to be empty.

\textbf{This is the deliverable.} A system that could not be wrong here
would not be doing science, and the demonstration that the architecture
\emph{can} destroy its own assertion is worth more than any assertion
that survives.

\begin{center}\rule{0.5\linewidth}{0.5pt}\end{center}

\subsubsection{4. The Cycle Turns: Revision and
Re-Jump}\label{the-cycle-turns-revision-and-re-jump}

Einstein's arrow from experiment back to Axioms does not mean
\emph{delete the axiom}. It means \textbf{revise it and go round again}.
A single disproof is not a cycle.

The revision is itself inductive, and it is constrained. The new axiom
must exclude the cases that killed the old one, and must still say
something. A revision that excluded each counterexample \textbf{by name}
would be unfalsifiable bookkeeping --- it would predict nothing about a
symbol not yet examined. The restriction must therefore be a
\textbf{structural predicate the graph can evaluate in advance}:

\begin{quote}
\textbf{Revised axiom.} Every function \emph{for which the graph records
at least one \texttt{takes\_arg} edge} takes at least one argument.

\textbf{Restriction.} The graph records a \texttt{takes\_arg} edge for
this symbol.
\end{quote}

The predicate is observable before any experiment is run, so the revised
axiom makes genuine predictions about unseen symbols.

\paragraph{4.1 Testing on held-out
symbols}\label{testing-on-held-out-symbols}

Testing the revision on the symbols that refuted it would be fitting a
claim to its own counterexamples. Six symbols not used in round 1 are
therefore held out. Symbols falling outside the restriction yield
\textbf{no prediction} --- which is a feature, since a claim that
predicts everything predicts nothing.

The revised axiom survives on the held-out symbols that fall inside its
scope, and is logged with provenance \texttt{induced}.

\textbf{Survival is not proof.} The claim is now harder to destroy, and
that is all. Induction confers no warrant beyond not-yet-refuted. What
the round demonstrates is \textbf{convergence}: a false universal became
a restricted claim that reality has not destroyed. Convergence is not
guaranteed --- that is a property of induction, not a defect of the
implementation.

\begin{center}\rule{0.5\linewidth}{0.5pt}\end{center}

\subsubsection{5. The Impossibility
Register}\label{the-impossibility-register}

\paragraph{5.1 The problem}\label{the-problem}

``X cannot be done'' is not one kind of claim, and treating it as one
has two opposite failure modes. A field can stall for a decade because
it mistook a narrow theorem for a general prohibition. A field can also
waste a decade chasing something genuinely closed by a conservation law.
Both errors come from discarding the \emph{warrant} of the claim and
keeping only its content.

The substrate therefore records, for every impossibility claim,
\textbf{how it is warranted} --- because warrant determines what a
successful counterexample means.

\paragraph{5.2 Four warrants}\label{four-warrants}

\begin{longtable}[]{@{}
  >{\raggedright\arraybackslash}p{(\linewidth - 4\tabcolsep) * \real{0.3333}}
  >{\raggedright\arraybackslash}p{(\linewidth - 4\tabcolsep) * \real{0.3333}}
  >{\raggedright\arraybackslash}p{(\linewidth - 4\tabcolsep) * \real{0.3333}}@{}}
\toprule\noalign{}
\begin{minipage}[b]{\linewidth}\raggedright
Warrant
\end{minipage} & \begin{minipage}[b]{\linewidth}\raggedright
Basis
\end{minipage} & \begin{minipage}[b]{\linewidth}\raggedright
What a counterexample does
\end{minipage} \\
\midrule\noalign{}
\endhead
\bottomrule\noalign{}
\endlastfoot
\textbf{ASSERTED} & someone said so & destroys it; nothing else
disturbed \\
\textbf{INDUCED} & no one has managed it yet & destroys it; the claim
was only a survey of failures \\
\textbf{DERIVED} & follows from a stated axiom & does \textbf{not}
refute the claim --- it refutes the \textbf{parent axiom}, and the
retraction must propagate upward \\
\textbf{THEOREM} & follows from a conservation law or proof & is a
\textbf{scope error until proven otherwise} \\
\end{longtable}

The asymmetry between the last two rows is the contribution.

For \textbf{DERIVED} claims, the counterexample lands somewhere other
than where it appears to. Retracting the impossibility while leaving the
axiom that entailed it is incoherent, and leaves every other consequence
of that axiom silently unsupported.

For \textbf{THEOREM} claims, the default reading of a counterexample is
that the new result \textbf{left the domain the theorem quantifies over}
--- not that a proof is broken. Before celebrating, the system must
answer: \emph{did this violate the law, or leave its scope?}

\paragraph{5.3 The canonical case}\label{the-canonical-case}

Minsky and Papert proved that a single-layer perceptron cannot compute
XOR. The theorem was never violated. Multilayer networks did not break
it; they \textbf{left its scope}, because the theorem quantifies over
single-layer perceptrons and a multilayer network is not one.

Read as a prohibition, the result froze a field. Read correctly, it was
a map to the exit: it named exactly the property --- single-layer ---
whose abandonment would be required.

\begin{quote}
\textbf{A good impossibility theorem does not close a field. It tells
you precisely where the bypass must lie.}
\end{quote}

\paragraph{5.4 The register, applied to this project's own
claims}\label{the-register-applied-to-this-projects-own-claims}

Self-application is deliberate. A system that records only \emph{other
people's} refuted impossibilities is not doing the thing this section
describes.

\textbf{(1) THEOREM --- escaped, intact.} \emph{``A bag of distinct
bytes cannot distinguish \texttt{metformin} from \texttt{metfromin}.''}
Rests on: a multiset is invariant under permutation. The theorem is true
and remains true. Bigram lanes are a different representation ---
ordered \emph{pairs}, not a cleverer weighting of the same multiset ---
so the sensory layer stopped being the object the theorem quantifies
over. Measured 7/7 on misspellings including transpositions.
\textbf{Scope escape. Not a violation.}

\textbf{(2) THEOREM --- escaped, intact.} \emph{``Global viscosity
modulation cannot break ground-state locking.''} Rests on: uniform
scaling preserves all amplitude ratios, and arg-max is invariant under
monotone transforms. The rank-preservation proof is untouched and still
correct. CSAT escaped by coupling per-cavity through
\(\tanh(T_i/T_{\text{crit}})\) --- it ceased to be \emph{global}. The
theorem pointed straight at the fix by naming the property that had to
go. \textbf{Scope escape. Not a violation.}

\textbf{(3) THEOREM --- stands, not escaped.} \emph{``A cognitive
architecture can eliminate the \(\log_2 N\) translation cost.''} Rests
on: the Translation Capacity Lemma (source coding). This is the honest
case. MRS relocates the bits from a label table into a coupling matrix,
changing their \emph{organisation} --- atomic to factored --- and
therefore what generalises. \textbf{The bit count is unchanged.} Any
claim to have removed the cost would be a scope error dressed as a
result.

\textbf{(4) INDUCED --- simply refuted.} \emph{``Every function takes at
least one argument.''} Rests on: 55\% observed frequency.
\texttt{tick()} and \texttt{reset\_field()} compile with zero arguments.
An induced universal needs no scope analysis; one counterexample ends
it.

\paragraph{5.5 The rule the substrate
applies}\label{the-rule-the-substrate-applies}

\begin{itemize}
\tightlist
\item
  \textbf{ASSERTED / INDUCED} --- one counterexample retracts the claim
  outright. Nothing else is disturbed.
\item
  \textbf{DERIVED} --- the counterexample does not touch the claim. It
  refutes the parent axiom, and retraction propagates to everything else
  that axiom supports.
\item
  \textbf{THEOREM} --- a counterexample is a scope error until proven
  otherwise. Two of the three theorem-claims above were escaped by
  leaving scope; \textbf{none was violated.}
\end{itemize}

A theorem that is \emph{genuinely} violated means the axioms behind it
are wrong. That is the largest event that can occur in this
architecture, and it has not occurred.

\begin{center}\rule{0.5\linewidth}{0.5pt}\end{center}

\subsubsection{6. Why This Is Not Reinforcement
Learning}\label{why-this-is-not-reinforcement-learning}

The distinction is worth stating precisely, because the loop
superficially resembles one.

\begin{longtable}[]{@{}
  >{\raggedright\arraybackslash}p{(\linewidth - 4\tabcolsep) * \real{0.3333}}
  >{\raggedright\arraybackslash}p{(\linewidth - 4\tabcolsep) * \real{0.3333}}
  >{\raggedright\arraybackslash}p{(\linewidth - 4\tabcolsep) * \real{0.3333}}@{}}
\toprule\noalign{}
\begin{minipage}[b]{\linewidth}\raggedright
\end{minipage} & \begin{minipage}[b]{\linewidth}\raggedright
Reinforcement learning
\end{minipage} & \begin{minipage}[b]{\linewidth}\raggedright
JAS
\end{minipage} \\
\midrule\noalign{}
\endhead
\bottomrule\noalign{}
\endlastfoot
Objective & expected return & none \\
Update & gradient on a value function or policy & retraction of a
claim \\
Signal & scalar reward & accept/reject verdict \\
Failure & lowers estimated value & \textbf{destroys an axiom} \\
Representation & learned parameters & typed graph edges \\
What accumulates & a policy & a record of what is false \\
\end{longtable}

There is no value function, no policy, no gradient, and no fitted
parameter. Nothing estimates a return. A refutation does not adjust a
weight --- it removes a statement from the system's commitments and logs
it as a region of hypothesis space now known to be empty.

The three neuromodulators the substrate carries (dopamine, cortisol,
serotonin) do move with outcome, but they are scalar multipliers on
physical constants --- drive, damping, and Hebbian rate --- with
homeostatic pull back to baseline. A run of refutations raises cortisol
and therefore damping, so the field literally propagates less far after
failure. This is affect, not credit assignment: nothing estimates
expected future reward, and no action is selected to maximise anything.

\begin{center}\rule{0.5\linewidth}{0.5pt}\end{center}

\subsubsection{7. Falsifiability of This Paper's Own
Claims}\label{falsifiability-of-this-papers-own-claims}

\begin{longtable}[]{@{}
  >{\raggedright\arraybackslash}p{(\linewidth - 2\tabcolsep) * \real{0.5000}}
  >{\raggedright\arraybackslash}p{(\linewidth - 2\tabcolsep) * \real{0.5000}}@{}}
\toprule\noalign{}
\begin{minipage}[b]{\linewidth}\raggedright
Claim
\end{minipage} & \begin{minipage}[b]{\linewidth}\raggedright
How to destroy it
\end{minipage} \\
\midrule\noalign{}
\endhead
\bottomrule\noalign{}
\endlastfoot
The jump is not entailed by the evidence & exhibit an inference rule
deriving the universal from 55\% frequency \\
Refutation is external & show the oracle shares state with the
substrate \\
The revision is not fitted to its counterexamples & show a held-out
symbol inside the restriction that refutes it \\
Theorem-claims (1) and (2) were escapes, not violations & exhibit a byte
\textbf{multiset} distinguishing the anagram pair, or a
\textbf{globally} uniform modulation breaking ground-state locking \\
Claim (3) stands & exhibit an architecture discriminating \(N\) outcomes
with \(< \log_2 N\) bits anywhere in the system \\
\end{longtable}

Each row is a concrete experiment, and the last two are the ones that
would matter most.

\begin{center}\rule{0.5\linewidth}{0.5pt}\end{center}

\subsubsection{8. Limitations}\label{limitations}

\textbf{One induction schema.} The system currently jumps over a single
relation (\texttt{takes\_arg}) with a single generalisation form. A
general mechanism would enumerate candidate universals across relation
types and rank them by refutability. Nothing in the architecture forbids
this; it is unwritten.

\textbf{Revision is hand-shaped.} The restriction predicate is chosen by
the implementer, not searched for. An automatic version would need a
principled bias toward restrictions that remain predictive on unseen
symbols --- and that bias is itself an inductive assumption which the
register would then have to record.

\textbf{The register is curated.} Its four entries were entered by a
human who noticed them. A system that detected its own impossibility
claims and classified their warrant automatically would be a stronger
result, and does not exist.

\textbf{One oracle.} The compiler settles questions about compilability
and nothing else. Domains without a cheap deterministic oracle --- most
empirical science --- do not inherit this loop for free.

\begin{center}\rule{0.5\linewidth}{0.5pt}\end{center}

\subsubsection{9. Conclusion}\label{conclusion}

The J-A-S cycle is implementable, and implementing it costs almost
nothing: induction is counting, deduction is graph traversal, and
refutation is an exit code. What it buys is a system that can hold a
claim, stake it on an outcome, and lose.

The impossibility register is the more portable contribution. Recording
the \emph{warrant} of a ``cannot'' --- rather than only its content ---
turns impossibility results from prohibitions into maps. Under this
discipline the question after an apparent breakthrough is not \emph{``is
the theorem broken?''} but \emph{``which quantified property did we
abandon?''} --- and in every case examined here, including two of the
project's own, the answer named the exit precisely.

The corresponding discipline for a claim that a system \emph{can} do
something is the one this lineage learned the hard way: re-derive it
from the running code before repeating it.

\begin{center}\rule{0.5\linewidth}{0.5pt}\end{center}

\subsubsection{References}\label{references}

\begin{itemize}
\tightlist
\item
  Einstein, A. Letter to Maurice Solovine, 7 May 1952. (The J-A-S
  diagram.)
\item
  Minsky, M. and Papert, S. \emph{Perceptrons}, MIT Press, 1969.
\item
  Popper, K. \emph{The Logic of Scientific Discovery}, 1959.
\item
  Mahir, I. \emph{From Artificial to Synthetic Intelligence: Principles
  of Constructive Structural Cognition.} §5.4c (experiment provenance),
  §5.4e (the cycle), §6.2 (lineage audit).
\item
  Mahir, I. \emph{Constructive Language Locking (CLL).} \texttt{cll.md}
  --- the category-error argument whose refutation pattern §5
  generalises.
\item
  Implementation: \texttt{src/jas.cpp} (cycle and register),
  \texttt{src/synth.cpp} (compile-verify oracle),
  \texttt{data/refuted.axioms} (refutation log).
\end{itemize}

\end{document}
